\chapter{Advanced File Handling: CSV and JSON in Python}

\section{Introduction: Real-Life Hook}

Students deal with data every day:  

\begin{itemize}
    \item Exam marks stored in spreadsheets  
    \item Contact lists saved in apps  
    \item Online data (like weather, scores) in JSON format  
\end{itemize}

Python allows us to \textbf{read, write, and manipulate structured data} using \textbf{CSV and JSON files}, making programs more useful in real life.

\section{Concept: What is CSV?}

\begin{definition}{CSV (Comma-Separated Values)}
A CSV file stores \textbf{tabular data} (like Excel sheets) in plain text, where each value is separated by a comma.
\end{definition}

\subsection{Reading a CSV File}

\begin{verbatim}
import csv

with open("marks.csv", "r") as file:
    reader = csv.reader(file)
    for row in reader:
        print(row)
\end{verbatim}

- Each \texttt{row} is a list of values  
- Use \texttt{csv.DictReader} to read rows as dictionaries:

\begin{verbatim}
with open("marks.csv", "r") as file:
    reader = csv.DictReader(file)
    for row in reader:
        print(row["Name"], row["Math"], row["Science"])
\end{verbatim}

\subsection{Writing to a CSV File}

\begin{verbatim}
import csv

data = [
    ["Name", "Math", "Science"],
    ["Amina", 80, 90],
    ["Rafi", 75, 85]
]

with open("marks.csv", "w", newline="") as file:
    writer = csv.writer(file)
    writer.writerows(data)
\end{verbatim}

- \texttt{writerows()} writes multiple rows at once  
- \texttt{newline=""} prevents extra blank lines  

\section{Concept: What is JSON?}

\begin{definition}{JSON (JavaScript Object Notation)}
A JSON file stores \textbf{data in key-value pairs}, making it easy to exchange information between programs and the web.
\end{definition}

\subsection{Reading a JSON File}

\begin{verbatim}
import json

with open("data.json", "r") as file:
    data = json.load(file)
    print(data)
\end{verbatim}

- \texttt{data} becomes a Python dictionary or list  

\subsection{Writing to a JSON File}

\begin{verbatim}
import json

data = {
    "students": [
        {"name": "Amina", "grade": 9},
        {"name": "Rafi", "grade": 8}
    ]
}

with open("data.json", "w") as file:
    json.dump(data, file, indent=4)
\end{verbatim}

- \texttt{indent=4} formats JSON neatly  

\section{Mini-Project: Student Marks Tracker}

Create a program to save and read student marks using CSV and JSON:

\begin{verbatim}
# CSV Version
import csv

students = [["Name", "Math", "Science"],
            ["Amina", 80, 90],
            ["Rafi", 75, 85]]

with open("marks.csv", "w", newline="") as file:
    writer = csv.writer(file)
    writer.writerows(students)

# Read CSV
with open("marks.csv", "r") as file:
    reader = csv.DictReader(file)
    for row in reader:
        print(row["Name"], row["Math"], row["Science"])

# JSON Version
import json

data = {"students": [{"name": "Amina", "math": 80, "science": 90},
                     {"name": "Rafi", "math": 75, "science": 85}]}

with open("marks.json", "w") as file:
    json.dump(data, file, indent=4)

# Read JSON
with open("marks.json", "r") as file:
    data = json.load(file)
    for student in data["students"]:
        print(student["name"], student["math"], student["science"])
\end{verbatim}

\section{Step-by-Step Diagram: CSV and JSON Workflow}

\begin{center}
\begin{tikzpicture}[node distance=1.5cm]
\node (start) [draw, ellipse, fill=green!20] {Start Program};
\node (write) [draw, rectangle, fill=yellow!20, below of=start] {Write Data to CSV / JSON};
\node (read) [draw, rectangle, fill=orange!20, below of=write] {Read Data from CSV / JSON};
\node (process) [draw, rectangle, fill=yellow!20, below of=read] {Process / Display Data};
\node (end) [draw, ellipse, fill=green!20, below of=process] {End Program};
\draw[->] (start) -- (write);
\draw[->] (write) -- (read);
\draw[->] (read) -- (process);
\draw[->] (process) -- (end);
\end{tikzpicture}
\end{center}

\section{Key Concepts}

\begin{itemize}
    \item CSV stores \textbf{tabular data}; JSON stores \textbf{structured key-value data}  
    \item Use \texttt{csv} and \texttt{json} modules to read/write files  
    \item CSV rows can be lists or dictionaries  
    \item JSON data maps naturally to Python dictionaries and lists  
    \item Useful for \textbf{data storage, exchange, and analysis}  
\end{itemize}

\section{End-of-Chapter Exercises}

\begin{exercise}{Very Short Questions}
\begin{enumerate}
    \item What does CSV stand for?  
    \item What does JSON stand for?  
    \item Name one Python module for CSV and one for JSON.  
\end{enumerate}
\end{exercise}

\begin{exercise}{Short Answer Questions}
\begin{enumerate}
    \item Write code to read a CSV file containing student marks and print each student’s total marks.  
    \item Write code to create a JSON file containing your favorite books and their authors.  
    \item Explain the difference between reading CSV with \texttt{csv.reader} and \texttt{csv.DictReader}.  
\end{enumerate}
\end{exercise}

\begin{exercise}{Activity-Based Questions}
\begin{enumerate}
    \item Create a CSV file for 5 students with Math and Science marks. Read the file and calculate the average for each student.  
    \item Create a JSON file storing employees’ names, positions, and salaries. Read it and display all employees earning more than 50,000.  
    \item Combine CSV and JSON: Read student data from CSV and save it into a JSON file.  
\end{enumerate}
\end{exercise}

\section{Teacher Notes}

\marginnote{
\textbf{Learning Objectives:}  
\begin{itemize}
    \item Students can read and write CSV and JSON files  
    \item Students understand mapping between CSV/JSON and Python data structures  
    \item Students can use structured files for real-life applications  
\end{itemize}

\textbf{Teaching Tips:}  
\begin{itemize}
    \item Demonstrate with small datasets first  
    \item Show CSV as spreadsheet analogy and JSON as key-value web data  
    \item Encourage students to combine CSV and JSON for projects  
\end{itemize}

\textbf{Assessment:}  
\begin{itemize}
    \item Check correct reading/writing of CSV and JSON  
    \item Observe correct mapping to Python dictionaries and lists  
\end{itemize}
}